\chapter{第八届大唐杯本科组省赛}
\fancyhead[L]{\color{H6}\kaishu\faIcon{atom}\;第八届大唐杯本科组省赛}
\fancyhead[R]{\color{H6}\kaishu\rightmark\,}

\date{2021年5月6日}{}{\href{https://github.com/xubenshan/dtcup}{\textbf{才疏学浅的小熊}}}
{}
{https://www.xiaohongshu.com/user/profile/6192219d0000000021028647}{小红书}
{https://space.bilibili.com/692546609}{B站}

\vspace{1em}
注：第八届省赛不分AB组，只有本科组。虽然比赛只有120道题，但每个人题目不一样，因此收集了144道题。

\section{单选题（共69小题）}



\begin{choice}{C}[]
    IP 地址分为哪两部分
    \begin{tasks}(1)
        \task 网络号和端口号
        \task 超网号和端口地址
        \task 网络号和主机号
        \task 超网号和子网号
    \end{tasks}
\end{choice}


\begin{choice}{B}[]
    5G mMTC 场景下，每平方公里至少支持多少台设备
    \begin{tasks}(4)
        \task 1000
        \task 100 万
        \task 1 万
        \task 10 万
    \end{tasks}
\end{choice}




\begin{choice}{A}[]
    下列哪项不属于 5G 网络逻辑架构关键指标
    \begin{tasks}(4)
        \task 高成本
        \task 低时延
        \task 低功率
        \task 高可靠
    \end{tasks}
\end{choice}




\begin{choice}{B}[]
    大规模天线形成阵列技术的英文简称是
    \begin{tasks}(4)
        \task mMTC
        \task Massive MIMO
        \task MIMO
        \task MU-MIMO
    \end{tasks}
\end{choice}


\begin{choice}{B}[]
    5G 核心网架构采用什么架构
    \begin{tasks}(4)
        \task 矩阵型
        \task 总线型
        \task 网状型
        \task 星型
    \end{tasks}
\end{choice}

\begin{choice}{A}[]
    5G NR 系统中，RRC INACTVE 状态，如果要恢复业务，则会出现哪条信令
    \begin{tasks}(2)
        \task RRC resume request
        \task RRC connection request
        \task RRC Connection reconfiguration
        \task RRC connection reject
    \end{tasks}
\end{choice}




\begin{choice}{C}[]
    5G NR 小区 SA 部署时，Initial DL BWP 的 BW、SCS 和 CP 由下面哪个 CORESET 定义
    \begin{tasks}(1)
        \task CORESET1
        \task CORESET2
        \task CORESET0
        \task CORESET3
    \end{tasks}
\end{choice}


\begin{choice}{A}[]
    5G NR 系统 RRC 重建成功流程，若 UE 上下文不能被恢复，基站会触发哪条信令
    \begin{tasks}(1)
        \task RRC Setup
        \task 不回任何消息
        \task RRC establishment
        \task RRC Reject
    \end{tasks}
\end{choice}


\begin{choice}{D}[]
    以下选项中哪个不是 5G 下行物理信道
    \begin{tasks}(4)
        \task PBCH
        \task PDSCH
        \task  PDCCH
        \task PCFICH
    \end{tasks}
\end{choice}


\begin{choice}{D}[]
    在 5G 无线网络规划中使用规划软件仿真时，修改以下哪项数值不需要重新计算路损
    \begin{tasks}(1)
        \task 天线型号
        \task 站点位置
        \task 天线下倾角
        \task 基站发射功率
    \end{tasks}
\end{choice}

\begin{choice}{D}[]
    对于大唐 5G 基站设备，HSCTD 板卡插在 0 槽位，则其登录的物理 IP 地址为
    \begin{tasks}(4)
        \task 172.27.245.92
        \task 172.27.245.100
        \task 172.27.45.250
        \task 172.27.245.91
    \end{tasks}
\end{choice}


\begin{choice}{D}[]
    对于大唐 64TRAAU，小区降质的定义是
    \begin{tasks}(1)
        \task AAU 有三分之一或者三分之一以上通道故障即可造成小区降质
        \task AAU 有二分之一或者二分之以上通道故障即可造成小区降质
        \task AAU 有 8 个以上通道故障即可造成小区降质
        \task AAU 有四分之一或者四分之一以上通道故障即可造成小区降质
    \end{tasks}
\end{choice}


\begin{choice}{B}[]
    同频小区重选邻区信息包含在哪个 SIB 中
    \begin{tasks}(2)
        \task SIB2
        \task SIB3
        \task SIB4
        \task  SIB5
    \end{tasks}
\end{choice}


\begin{choice}{B}[]
    如下5G的关键技术中，哪个不可以提升网络频谱效率
    \begin{tasks}(2)
        \task 新信道编码
        \task 灵活子载波带宽
        \task 256QAM 高阶调制
        \task Massive \;MIMO
    \end{tasks}
\end{choice}

\begin{choice}{D}[]
    车联网是实现自动驾驶的必要条件，以下不属于车联网优点的为
    \begin{tasks}(1)
        \task 提升安全性
        \task 可提供丰富的网联应用
        \task 提升交通效率
        \task 可增加感知范围，但感知成本会相应提升
    \end{tasks}
\end{choice}

\begin{choice}{C}[]
    5G NR 中 UE 如果在 slot n 收到 timing advance command，那么应该在什么时刻调整上行传输timing
    \begin{tasks}(2)
        \task slot\;n+2
        \task slot\;n+4
        \task slot\;n+6
        \task slot\;n+8
    \end{tasks}
\end{choice}



\begin{choice}{C}[]
    5G NSA 网络场景下，哪一层的统计更能准确反映用户在 5G NR 获得的下行速率
    \begin{tasks}(4)
        \task 应用层
        \task RRC 层
        \task PDCP 层
        \task RLC 层
    \end{tasks}
\end{choice}

\begin{choice}{D}[]
    在 5G 网络架构中，以下选项哪一项是 AMF 的功能
    \begin{tasks}(4)
        \task 下行数据的通知
        \task 漫游功能
        \task 会话的建立修改删除
        \task 注册管理
    \end{tasks}
\end{choice}


\begin{choice}{D}[]
    在 5G 网络中，FTP 文件下载业务属于如下那一种业务场景
    \begin{tasks}(2)
        \task 以上全部
        \task 海量机器类(mMTC)
        \task 超可靠低时延( URLLC)
        \task 增强型移动宽带(eMBB)
    \end{tasks}
\end{choice}


\begin{choice}{B}[]
    单模光纤是什么类型光纤
    \begin{tasks}(2)
        \task 渐变型
        \task 阶跃型
        \task 梯度型
        \task 自聚焦型
    \end{tasks}
\end{choice}



\begin{choice}{B}[]
    5G 中从 BBU 到 AAU 需要保证多少Gbps 的传输带宽
    \begin{tasks}(2)
        \task 15
        \task 25
        \task 10
        \task 5
    \end{tasks}
\end{choice}




\begin{choice}{C}[]
    5G NR 中若 PRACH 为短序列，在中低速场景，零相关配置为 11，对应的 Ncs 为 23，单小区需要多少个逻辑根序列
    \begin{tasks}(4)
        \task 10
        \task 12
        \task 11
        \task 9
    \end{tasks}
\end{choice}

\begin{choice}{A}[]
    5G NR 系统中，100M 带宽组网，30kHz 子载波间隔，单小区最大支持的 PRB 为
    \begin{tasks}(2)
        \task 273PRB
        \task 175PRB
        \task 100PRB
        \task 200PRB
    \end{tasks}
\end{choice}

\begin{choice}{B}[]
    5G 通信系统中，gNB 向 INACTIVE 态的 UE 发送 Paging 消息时，使用的寻呼只为
    \begin{tasks}(2)
        \task P-RNTI
        \task I-RNTI
        \task S-TMSI
        \task P-RNTI
    \end{tasks}
\end{choice}


\begin{choice}{D}[]
    第一个支持5G-V2X 标准的 3GPP 协议版本为
    \begin{tasks}(1)
        \task  R14
        \task  R13
        \task R15
        \task R16
    \end{tasks}
\end{choice}


\begin{choice}{C}[]
    5G NR 构架中数据传输的方式不包括
    \begin{tasks}(2)
        \task 回传
        \task 中传
        \task 后传
        \task 前传
    \end{tasks}
\end{choice}



\begin{choice}{C}[]
    5G NR 系统中，物理小区组 ID 标识范围为
    \begin{tasks}(2)
        \task 0-2
        \task 0-3
        \task 0-335
        \task 0-503
    \end{tasks}
\end{choice}




\begin{choice}{D}[]
    5G NR系统中，关于随机接入的目的描述错误的是
    \begin{tasks}(2)
        \task 获取C-RNTI，基站识别UE的标识
        \task 获取MSG3的资源
        \task 实现UE和gNB之间的上行同步
        \task 实现UE和gNB之间的下行同步
    \end{tasks}
\end{choice}


\begin{choice}{B}[]
    5G NR系统中，SPS配置使用如下哪种RNTI功扰PDCCH进行资源调度
    \begin{tasks}(2)
        \task RA-RANTI
        \task CS-RNTI
        \task  C-RNTI
        \task SPS C-RNTI
    \end{tasks}
\end{choice}


\begin{choice}{B}[]
    关于大唐5GRAN设备（EMB6116）技术特点描述正确的是
    \begin{tasks}(2)
        \task AAU接口仅支持25G光接口
        \task 支持NR新型帧结构、新参数集、新型编码
        \task 基带单板300MHz处理能力
        \task 仅支持GPS、BD时钟获取方式
    \end{tasks}
\end{choice}

\begin{choice}{D}[]
    EMB6116未使用的槽位必须加装空面板，否则影响散热。下列（\quad）槽位为空时需安装带导风板的空面板。
    \begin{tasks}(2)
        \task SLOT1
        \task SLOT3
        \task SLOT2
        \task SLOT6
    \end{tasks}
\end{choice}


\begin{choice}{D}[]
    5G NR系统中，MIB一个最重要作用是通知UE如何获取（\quad）消息
    \begin{tasks}(2)
        \task Msg2
        \task Msg1
        \task SIB2
        \task SIB1
    \end{tasks}
\end{choice}


\begin{choice}{B}[]
    大唐5G基站设备-EMB6116，HBPOD板卡不具备的功能
    \begin{tasks}(1)
        \task SDAP
        \task 与BBU内部各板卡之间的业务、信令交换处理功能
        \task L2处理功能
        \task 物理层处理功能
    \end{tasks}
\end{choice}


\begin{choice}{B}[]
    5G NR通信系统中，CORESET在时域上最多可以使用几个OFDM符号
    \begin{tasks}(2)
        \task 4
        \task 3
        \task 2
        \task 1
    \end{tasks}
\end{choice}



\begin{choice}{A}[]
    大唐5G基站设备，BBU接地线采用多少平方黄绿接地线？
    \begin{tasks}(2)
        \task 16
        \task 25
        \task 6
        \task 35
    \end{tasks}
\end{choice}

\begin{choice}{A}[]
    目前，我国5G商用网采用的NSA组网部署方式为哪一种

    \begin{tasks}(2)
        \task option3x
        \task option7
        \task option7x
        \task option3
    \end{tasks}
\end{choice}


\begin{choice}{A}[]
    目前工信部已经发放车联网直连通信频率，在中国车联网直连通信使用频率为
    \begin{tasks}(2)
        \task 5905-5925MHz
        \task 5855-5925MHz
        \task 5850-5925MHz
        \task 5770-5850MHz
    \end{tasks}
\end{choice}


\begin{choice}{B}[]
    5G技术中，用于提升接入用户数的技术是
    \begin{tasks}(2)
        \task Massive CA
        \task Massive MIMO
        \task NOMA
        \task CDMA
    \end{tasks}
\end{choice}


\begin{choice}{D}[]
    5G NR通信系统中，小区选择的准则是
    \begin{tasks}(2)
        \task Srxlev>0或Squal<0
        \task Srxlev>0或Squal>0
        \task Srxlev>0且Squal<0
        \task Srxlev>0且Squal>0
    \end{tasks}
\end{choice}


\begin{choice}{D}[]
    5G控制信道资源调度的基本单位CCE与REG之间的正确关系为
    \begin{tasks}(2)
        \task 1CCE = 7REG
        \task 1CCE = 8REG
        \task 1CCE = 9REG
        \task 1CCE = 6REG
    \end{tasks}
\end{choice}

\begin{choice}{C}[]
    5G网络架构中，CU/DU分离场景下，回传一般指哪两个之间的数据传送
    \begin{tasks}(2)
        \task AAU-DU
        \task AAU-BBU
        \task CU-5GC
        \task DU-CU
    \end{tasks}
\end{choice}

\begin{choice}{A}[]
    在LTE-V2X中，基于PC5接口，基站集中调度分配模式是
    \begin{tasks}(2)
        \task mod3
        \task mod2
        \task mod5
        \task mod4
    \end{tasks}
\end{choice}





\begin{choice}{A}[]
    5G NR系统中，PDCCH信道采用Polar码信道编码方式，调制方式为
    \begin{tasks}(2)
        \task QPSK
        \task 64QAM
        \task 16QAM
        \task BPSK
    \end{tasks}
\end{choice}


\begin{choice}{C}[]
    5G NR系统中，基于SSB的NR同频测量在measconfig里最多可以配置多少个SMTC窗口
    \begin{tasks}(2)
        \task 4
        \task 1
        \task 2
        \task 3
    \end{tasks}
\end{choice}

\begin{choice}{D}[]
    5G NR系统中基于基站间Xn链路切换，源基站侧向目标基站发送第一条信令为
    \begin{tasks}(2)
        \task Handover Request Acknowledge
        \task SN Status Transfer
        \task Handover Required
        \task Handover Request
    \end{tasks}
\end{choice}


\begin{choice}{D}[]
    对于大唐5G设备，单模集束光纤在BBU与AAU路由长度超过多少米使用。
    \begin{tasks}(2)
        \task 80
        \task 60
        \task 40
        \task 100
    \end{tasks}
\end{choice}


\begin{choice}{D}[]
    5G NR支持几种长PRACH格式。
    \begin{tasks}(4)
        \task 3
        \task 8
        \task 9
        \task 4
    \end{tasks}
\end{choice}


\begin{choice}{D}[]
    5G NR帧结构中当采用常规CP时，一个时隙中固定有多少个符号
    \begin{tasks}(2)
        \task 12
        \task 20
        \task 6
        \task 14
    \end{tasks}
\end{choice}


\begin{choice}{B}[]
    在超密集组网的场景中，什么是解决边界效应的关键
    \begin{tasks}(2)
        \task 软扇区
        \task 小区虚拟化
        \task 全频谱接入
        \task 回传技术
    \end{tasks}
\end{choice}


\begin{choice}{B}[]

    5G通信系统上行频谱效率的目标是多少bps/hz
    \begin{tasks}(2)
        \task 20
        \task 15
        \task 30
        \task 10
    \end{tasks}
\end{choice}


\begin{choice}{A}[]
    对于大唐5G基站设备维护时，当传输状态正常时，HSCTD指示灯GE、Ofp1应显示为
    \begin{tasks}(2)
        \task 常亮
        \task 慢闪
        \task 不亮
        \task 快闪
    \end{tasks}
\end{choice}


\begin{choice}{D}[]
    NR系统中，下列哪项不属于5G无线技术
    \begin{tasks}(2)
        \task 新型多址
        \task 大规模天线阵列
        \task 超密集组网
        \task 网络功能虚拟化
    \end{tasks}
\end{choice}

\begin{choice}{B}[]
    下列哪个因素不影响AAU正常接入
    \begin{tasks}(2)
        \task 对应的HBPOD板卡状态为未初始化状态
        \task 修改了小区物理ID列表
        \task DCPD给AAU供电的开关未闭合
        \task 对应的光模块插反了
    \end{tasks}
\end{choice}

\begin{choice}{C}[]
    在5G NR帧结构中，当SCS子载波间隔为30kHz时一个子帧包含的符号数为
    \begin{tasks}(2)
        \task 56个
        \task 48个
        \task 28个
        \task 14个
    \end{tasks}
\end{choice}


\begin{choice}{C}[]
    对于大唐5G设备EMB6116，HDPSD输出电压为多少伏特

    \begin{tasks}(2)
        \task 24
        \task 5
        \task 12
        \task 48
    \end{tasks}
\end{choice}


\begin{choice}{A}[]
    对于大唐5G基站设备，AAU进行近端直连升级时，需要登录的IP地址为
    \begin{tasks}(2)
        \task 172.27.245.100
        \task 172.27.45.250
        \task 172.27.245.91
        \task 172.27.245.92
    \end{tasks}
\end{choice}


\begin{choice}{C}[]
    5G SA场景下，本地cell ID的取值范围是
    \begin{tasks}(2)
        \task 0-127
        \task 0-63
        \task 0-11
        \task 0-255
    \end{tasks}
\end{choice}

\begin{choice}{B}[]
    在5G网络架构中，我们所提到的软件和硬件解耦主要通过下面哪种技术实现
    \begin{tasks}(2)
        \task SDN
        \task NFV
        \task 切片技术
        \task 移动边缘计算
    \end{tasks}
\end{choice}




\begin{choice}{B}[]
    5G NR C-V2X中，车与基站之间直连通信接口是
    \begin{tasks}(2)
        \task Xn
        \task Uu
        \task Ng
        \task PC5
    \end{tasks}
\end{choice}


\begin{choice}{B}[]
    5G NR 大唐5G基站安装时，2×16平电源线长期允许载流量是
    \begin{tasks}(2)
        \task 45A
        \task 76A
        \task 80A
        \task 100A
    \end{tasks}
\end{choice}


\begin{choice}{C}[]
    5G NR系统中，5G基站和核心网通过哪个接口相连
    \begin{tasks}(2)
        \task S1
        \task X2
        \task NG
        \task Xn
    \end{tasks}
\end{choice}


\begin{choice}{B}[]
    5G NR系统中PCI取值范围为
    \begin{tasks}(2)
        \task 1-504
        \task 0-1007
        \task 1-1008
        \task 0-503
    \end{tasks}
\end{choice}

\begin{choice}{C}[]
    5G NR系统中，对于SA场景，UE ip address allocation 功能由哪个功能模块实现
    \begin{tasks}(2)
        \task UPF
        \task MME
        \task SMF
        \task AMF
    \end{tasks}
\end{choice}


\begin{choice}{B}[]
    3GPP协议针对5G第一个商用版本为
    \begin{tasks}(4)
        \task R14
        \task R15
        \task R17
        \task R16
    \end{tasks}
\end{choice}

\begin{choice}{B}[]
    对于大唐5G基站设备，GPS天线的工作电压为
    \begin{tasks}(2)
        \task 直流-12V
        \task 直流5V
        \task 交流220V
        \task 直流-48V
    \end{tasks}
\end{choice}


\begin{choice}{B}[]
    在通信系统中，当无线电波遇到远大于波长的障碍物表面时，会产生什么现象
    \begin{tasks}(2)
        \task 散射
        \task 反射
        \task 直射
        \task 绕射
    \end{tasks}
\end{choice}


\begin{choice}{B}[]
    5G NR系统中，下列哪个选项中的PRACH配置周期不正确
    \begin{tasks}(4)
        \task 10ms
        \task 60ms
        \task 40ms
        \task 160ms
    \end{tasks}
\end{choice}


\begin{choice}{C}[]
    NetworkSlicing技术中，当用户同时接入多个切片时，其可以分为两种架构，其一为独立架构，另外一种是
    \begin{tasks}(2)
        \task 子网架构
        \task 非独立架构
        \task 共享架构
        \task 物理架构
    \end{tasks}
\end{choice}


\begin{choice}{A}[]
    对于大唐5G基站设备，当时钟同步时，至少需要锁定几颗卫星
    \begin{tasks}(4)
        \task 3
        \task 4
        \task 5
        \task 6
    \end{tasks}
\end{choice}




\section{多选题（共51小题）}

\begin{choice}{\;BCD\;}[]
    5G NR 中，包括哪几种栅格
    \begin{tasks}(2)
        \task Measure Raster
        \task Channel Raster
        \task SS Raster
        \task  Global Raster
    \end{tasks}
\end{choice}

\begin{choice}{\;ACD\;}[]
    5G网络部署Massive\;MIMO主要用于解决哪些问题
    \begin{tasks}(2)
        \task 深度覆盖问题
        \task 穿透损耗问题
        \task 天面不足问题
        \task 大容量问题
    \end{tasks}
\end{choice}

\begin{choice}{\;ABD\;}[]
    以下哪些属于中国运营商的5G频段
    \begin{tasks}(2)
        \task 3500-3600MHz
        \task 3400-3500MHz
        \task 1880-1900MHz
        \task 4800-4900MHz
    \end{tasks}
\end{choice}

\begin{choice}{\;BC\;}[]
    5G NR 系统中针对系统消息部分采用了不同的方式，其中最小系统消息包括
    \begin{tasks}(4)
        \task SIB2
        \task SIB1
        \task MIB
        \task SIB3
    \end{tasks}
\end{choice}


\begin{choice}{\;BD\;}[]
    在通信系统中，一般情况下，时延包含的内容是
    \begin{tasks}(2)
        \task 反应时延
        \task 处理时延
        \task 计算时延
        \task 传输时延
    \end{tasks}
\end{choice}

\begin{choice}{\;BD\;}[]
    与 4G 网络相比，5G 网络架构更加注重软件技术，包括了网络的
    \begin{tasks}(2)
        \task 蜂窝化
        \task 云化
        \task 集中化
        \task 虚拟化
    \end{tasks}
\end{choice}

\begin{choice}{\;CD\;}[]
    5G NR 系统中，需要进行功率控制的信道/信号包括
    \begin{tasks}(2)
        \task CRS
        \task PHICH
        \task SRS
        \task PUSCH
    \end{tasks}
\end{choice}


\begin{choice}{\;BCD\;}[]
    5G D2D 通信优势是
    \begin{tasks}(2)
        \task 无线 M2M 功能
        \task 拓展网络覆盖范围
        \task 短距离通信可频谱资源复用
        \task 终端近距离通信，高速率低时延低功耗
    \end{tasks}
\end{choice}


\begin{choice}{\;BCD\;}[]
    5G NR 中，RRC 状态包括哪些
    \begin{tasks}(2)
        \task URA CONNECTED
        \task RRC INACTIVE
        \task RRC IDLE
        \task RRC CONNECTED
    \end{tasks}
\end{choice}


\begin{choice}{\;CD\;}[]
    下列数据类型，哪些属于 5G 物联网
    \begin{tasks}(2)
        \task 传输类数据
        \task 交互类数据
        \task 控制类数据
        \task 采集类数据
    \end{tasks}
\end{choice}

\begin{choice}{\;ABCD\;}[]
    5G NR 系统中，以下哪些原因会导致 RRC 连接重建
    \begin{tasks}(2)
        \task 完保校验失败
        \task 重配置失败
        \task 检测到 RLF
        \task 切换失败
    \end{tasks}
\end{choice}

\begin{choice}{\;BC\;}[]
    5G 的架构哪个选项是独立架构
    \begin{tasks}(4)
        \task 4
        \task 5
        \task 2
        \task 3
    \end{tasks}
\end{choice}

\begin{choice}{\;BD\;}[]
    根据自由空间传播模型可知路径损耗与以下哪些因素有
    \begin{tasks}(2)
        \task 发射功率
        \task 传播距离
        \task 天线增益
        \task 工作频率
    \end{tasks}
\end{choice}

\begin{choice}{\;ABD\;}[]
    以下哪些选项属于 NR-V2X 可支持传输模式
    \begin{tasks}(4)
        \task 广播
        \task 单播
        \task 分播
        \task 组播
    \end{tasks}
\end{choice}

\begin{choice}{\;ACD\;}[]
    以下选项关于 5G 帧结构描述正确的是
    \begin{tasks}(1)
        \task 每个subframe的slot个数与scs相关
        \task 每个 subframe 的 slot 个数固定
        \task  slot及符号长度与scs有关
        \task 特殊子帧位置灵活配置
    \end{tasks}
\end{choice}

\begin{choice}{\;ACD\;}[]
    在5G所支持的SCS中，能被所有物理信道使用的SCS是
    \begin{tasks}(4)
        \task 15KHz
        \task 60KHz
        \task 30KHz
        \task 120KHz
    \end{tasks}
\end{choice}

\begin{choice}{\;ACD\;}[]
    5G R15版本中，信道状态信息CSI可能包括以下选项中的哪些内容
    \begin{tasks}(4)
        \task  CQI
        \task SINR
        \task PMI
        \task RI
    \end{tasks}
\end{choice}

\begin{choice}{\;ABCD\;}[]
    5G 的网络架构描述正确的是
    \begin{tasks}(1)
        \task 控制平面引入新的基于服务的设计思路
        \task 分布式的用户面功能 UPF，无需汇聚点，支持多会话
        \task 核心网抽取会话管理功能，形成独立模块 SMF
        \task 架构设计以网元功能为单位，支持网络切片
    \end{tasks}
\end{choice}

\begin{choice}{\;BC\;}[]
    关于 5G NR PCI 的描述，正确的是
    \begin{tasks}(1)
        \task 相邻小区间必须保证 PCI 模 3 不同
        \task 相邻小区间 PCI 模 30 可以相同
        \task 相邻小区间 PCl 模 3 可以相同
        \task 相邻小区间必须保证 PCI 模 30 不同
    \end{tasks}
\end{choice}

\begin{choice}{\;AC\;}[]
    在移动通信中对上行和下行概念理解正确的是
    \begin{tasks}(2)
        \task 上行代表从终端至基站
        \task 下行代表从终端至基站
        \task 下行代表从基站至终端
        \task 上行代表从基站至终端
    \end{tasks}
\end{choice}



\begin{choice}{\;CD\;}[]
    在 5G NR 中 SS Block burst 的时长为 5ms，5ms 内包含的 SS Block 个数与频段相关的个数定为 L，以下说法哪些正确
    \begin{tasks}(2)
        \task < 3GHz，L= 12
        \task < 3GHz，L=6
        \task 3GHz-6GHz，L=8
        \task 6GHz-52.6GHz，L=64
    \end{tasks}
\end{choice}

\begin{choice}{\;ABCD\;}[]
    NR 系统中PCF 的网元功能为
    \begin{tasks}(1)
        \task 实现前端（PCFFE），以访问统一数据存储库（UDR）
        \task PCF 访问与 PCF 相同的 PLMN 中的 UDR
        \task 提供控制面功能策略规则
        \task 支持统一的政策架构来管理网络行为
    \end{tasks}
\end{choice}


\begin{choice}{\;ABD\;}[]
    大唐 5G 设备维护过程中，以下哪些选项可能造成 RRU 和 BBU 之间的光纤连接异常
    \begin{tasks}(2)
        \task 光纤头子受污染
        \task 光纤长度过长
        \task 光纤走纤时曲率半径过大
        \task 光模块类型不匹配
    \end{tasks}
\end{choice}


\begin{choice}{\;AC\;}[]
    大唐 5G 基站开通后，关于板卡状态，说法正确的是
    \begin{tasks}(2)
        \task 管理状态应为解锁定状态
        \task 管理状态应为锁定状态
        \task 板卡过程状态应为初始化结束状态
        \task 板卡过程状态为未初始化状态
    \end{tasks}
\end{choice}


\begin{choice}{\;ABCD\;}[]
    EMB6116 支持的时钟源同步方式包括
    \begin{tasks}(2)
        \task GPS
        \task 1588 v2
        \task BD(北斗)
        \task 级联时钟
    \end{tasks}
\end{choice}


\begin{choice}{\;AC\;}[]
    在 5G 移动通信中影响信号传播的三种最基本的电磁波传播机制为
    \begin{tasks}(2)
        \task 反射
        \task 透射
        \task 散射
        \task 折射
    \end{tasks}
\end{choice}


\begin{choice}{\;CD\;}[]
    5G A3 事件配置中，以下哪些配置方法可以抑制切换
    \begin{tasks}(2)
        \task 减小源小区到目标小区的 CIO
        \task 减小 A3 上报触发时间
        \task 增加 A3 上报的 RSRP 偏移
        \task 增大 A3 上报间隔
    \end{tasks}
\end{choice}


\begin{choice}{\;ABC\;}[]
    5G NR 中加密算法包括
    \begin{tasks}(4)
        \task nea2
        \task nea3
        \task nea1
        \task nea4
    \end{tasks}
\end{choice}


\begin{choice}{\;AD\;}[]
    大唐 EMB6116 BBU 未使用的槽位必须加装空面板。下列哪些槽位为空时，需安装带导风板的空面板
    \begin{tasks}(2)
        \task SLOT6
        \task SLOT3
        \task SLOT2
        \task SLOT7
    \end{tasks}
\end{choice}


\begin{choice}{\;CD\;}[]
    UE 在上报 NR 邻区的 MR 中，会携带 NR 小区的哪些信息
    \begin{tasks}(2)
        \task  PathLoss
        \task TA
        \task PCI
        \task RSRP
    \end{tasks}
\end{choice}


\begin{choice}{\;BCD\;}[]
    以下选项中属于5G NR系统上行物理信道的是
    \begin{tasks}(2)
        \task  PDSCH
        \task PRACH
        \task PUCCH
        \task PUSCH
    \end{tasks}
\end{choice}

\begin{choice}{\;AC\;}[]
    以下选项中，属于 NR−V2X 资源分配模式的
    \begin{tasks}(2)
        \task Mode2
        \task Mode3
        \task Mode1
        \task Mode4
    \end{tasks}
\end{choice}

\begin{choice}{\;ABC\;}[]
    5G SA组网下，NR小区全局标识符NCGI由哪些组成
    \begin{tasks}(2)
        \task  gNB ID
        \task PLMN
        \task Cell ID
        \task AMF ID
    \end{tasks}
\end{choice}


\begin{choice}{\;ABCD\;}[]
    5G NR系统中，关于RRM测量配置描述正确的是
    \begin{tasks}(1)
        \task 测量事件支持A类、B类事件
        \task CSI-RS（SSB）触发的上报中不携带SSB（CSI-RS）的测量结果
        \task 测量量包括：RSRP/RSRQ/SINR
        \task 可配置测量事件基于SSB或CSI-RS
    \end{tasks}
\end{choice}



\begin{choice}{\;ABCD\;}[]
    对于大唐5G产品，以下选项关于AAU的描述正确的是
    \begin{tasks}(1)
        \task AAU集成了RRU和天线两个模块
        \task AAU具有简化天面、安装方便、加快建网的优点跃型
        \task AAU架构更有利于天线校准精度，减少由于线缆连接而造成的不可控因素，获得更好的波束赋型性能
        \task AAU的全称为ActiveAntennaUnit
    \end{tasks}
\end{choice}


\begin{choice}{\;AD\;}[]
    2020年，5G设备NG接口支持的速率包括
    \begin{tasks}(4)
        \task 25G
        \task 100G
        \task 50G
        \task 10G
    \end{tasks}
\end{choice}



\begin{choice}{\;AD\;}[]
    进行5G基站安装时，关于TDAU5164N78安装最小空间要求（单位：mm），以下说法正确的是
    \begin{tasks}(1)
        \task AAU 顶部应预留 300mm 布线和维护空间
        \task AAU左侧应预留350mm布线和维护空间
        \task AAU右侧应预留350mm布线和维护空间
        \task AAU底部应预留600mm布线空间，为方便维护建议底部距地面至少1200mm
    \end{tasks}
\end{choice}



\begin{choice}{\;ABC\;}[]
    5G NR系统中，关于PDCCH以下说法正确的是

    \begin{tasks}(1)
        \task CORESET在频域包含一组PRBS，最小粒度为6RB
        \task 控制信道由CCE聚合而成
        \task 在时域长度为1-3个 OFDM符号，在时隙中开始位置可配置
        \task CORESET内，CCE-to-REG的映射关系，仅支持非交织
    \end{tasks}
\end{choice}



\begin{choice}{\;ACD\;}[]
    大唐5G设备安装时，以下BBU竖装规范要求，描述正确的是
    \begin{tasks}(1)
        \task BBU竖装时，前方需留有800mm的维护通道
        \task BBU竖装时，机柜后部不支持靠墙安装，与墙的可距离800mm
        \task BBU竖装支持并柜或并架安装；
        \task BBU竖装双排或多排机柜安装时，机柜前后间隔≥800mm，同时需要同风向安装；
    \end{tasks}
\end{choice}


\begin{choice}{\;BCD\;}[]
    针对大唐5G产品，进行巡检工作时，常用的工具有哪些

    \begin{tasks}(2)
        \task OSP
        \task LMT
        \task NodeBSpider
        \task uem5000
    \end{tasks}
\end{choice}


\begin{choice}{\;ABCD\;}[]
    对于大唐5G设备，AAU版本升级前必须的检查项有哪些？
    \begin{tasks}(2)
        \task 先保证链路正常，光口光纤同步
        \task AAU能够正常接入
        \task AAU与BBU连接是否正常
        \task OSP上相应处理器能连接上
    \end{tasks}
\end{choice}


\begin{choice}{\;ABCD\;}[]
    属于5G超密集组网应用场景有
    \begin{tasks}(2)
        \task 商业中心
        \task 密集住宅区
        \task 大型活动场馆
        \task 交通枢纽
    \end{tasks}
\end{choice}


\begin{choice}{\;ABC\;}[]
    在5G NR中PRACH规划包含哪些
    \begin{tasks}(2)
        \task Preamble格式规划
        \task 根序列索引
        \task PRACH配置索引
        \task SSS配置规划
    \end{tasks}
\end{choice}


\begin{choice}{\;AD\;}[]
    大唐5G设备安装时，如果发现现场空开规格不满足要求，以下做法正确的是
    \begin{tasks}(2)
        \task 暂时不接电，先进行设备安装其他作业
        \task 艺高人胆大”继续接电作业
        \task 直接对空开进行改造使其符合开通要求
        \task 上报客户现场存在的问题
    \end{tasks}
\end{choice}



\begin{choice}{\;ACD\;}[]
    大唐5G基站产品EMB6216，HBPOF可以插在哪些槽位
    \begin{tasks}(4)
        \task  3、7
        \task 0、4
        \task 2、6
        \task 1、5
    \end{tasks}
\end{choice}


\begin{choice}{\;BC\;}[]
    5G网络巡检使用LMT时，随LMT同步打开的两个软件分别是
    \begin{tasks}(2)
        \task LMT
        \task LmtAgent
        \task FTP server
        \task FTP client
    \end{tasks}
\end{choice}

\begin{choice}{\;CD\;}[]
    现阶段，5G网络部署方式有两种，分别是
    \begin{tasks}(2)
        \task 联合部署
        \task 半独立部署
        \task 非独立部署
        \task 独立部署
    \end{tasks}
\end{choice}

\begin{choice}{\;ACD\;}[]
    5G NR系统中，MIB消息中包含以下哪些内容
    \begin{tasks}(2)
        \task SIB1的PDCCH CORESET配置
        \task PHICH配置
        \task 半帧指示
        \task 系统帧号
    \end{tasks}
\end{choice}

\begin{choice}{\;ABCD\;}[]
    5G SA场景下，以下属于NG-RAN的功能是
    \begin{tasks}(2)
        \task 无线资源管理、会话管理
        \task 包头压缩，加密和完整性保护
        \task 调度，传输寻呼、系统广播消息
        \task QoS流映射
    \end{tasks}
\end{choice}

\begin{choice}{\;BD\;}[]
    5G系统中，以下哪些应用属于URLLC场景
    \begin{tasks}(2)
        \task 视频会议
        \task 自动驾驶
        \task 水电煤等基础设施计费
        \task 设备远程控制
    \end{tasks}
\end{choice}

\begin{choice}{\;AD\;}[]
    下面有关网络切片说法正确的是
    \begin{tasks}(1)
        \task 每个切片之间是逻辑上隔离
        \task 在一个硬件基础设施切分出多个物理的端到端网络
        \task 每个切片之间是物理上隔离
        \task 在一个硬件基础设施切分出多个虚拟的端到端网络
    \end{tasks}
\end{choice}




\section{判断题（共24小题）}

\begin{choice}{\;正确\;}[]
    5G NR 系统中下行初始 BWP 占用的连续 PRB 资源，可以与 CORESET 的 Type0 PDCCH CSS 集合相同。
\end{choice}


\begin{choice}{\;错误\;}[]
    在 5G NR 中，NZC 代表随机接入前导序列长度，取值可能为 838 和 138。
\end{choice}

\begin{choice}{\;正确\;}[]
    MEC 是 Mobile Edge Computing，即移动边缘计算。
\end{choice}

\begin{choice}{\;正确\;}[]
    大唐 5G 网络版本默认的时钟源是 GPS。
\end{choice}

\begin{choice}{\;正确\;}[]
    5G 中的子帧和 4G 中的子帧时长都是 1ms。
\end{choice}

\begin{choice}{\;正确\;}[]
    大唐 5G 基站设备EMB6116，HBPOD 板卡 RUN 灯快闪表示板卡正在固件升级。

\end{choice}


\begin{choice}{\;正确\;}[]
    借助于人，车，路，云平台之间的全方位连接和高效信息交互，C-V2X 目前正从信息服务类应用向交通安全和效率类应用发展，并将逐步向支持实现自动驾驶的协同服务类应用演进。

\end{choice}

\begin{choice}{\;正确\;}[]
    5G 的 AAU 设备包括天线，所以不用额外部署天线。

\end{choice}

\begin{choice}{\;正确\;}[]

    5G NR 系统中，UE 可通过 SIB3 获取同频小区重选的相关消息。
\end{choice}


\begin{choice}{\;错误\;}[]
    5G NR 系统中，UE 在 RRC IDLE 态和 RRC INACTIVE 态的寻呼标识为 P-RNTI。

\end{choice}

\begin{choice}{\;错误\;}[]
    大唐 5G 产品 EMB6116 机框仅支持 8/9/10 槽位插装 HBPOF 单板。

\end{choice}


\begin{choice}{\;正确\;}[]
    5G SA 场景下，业务转发的配置不属于 UPF 的功能。

\end{choice}


\begin{choice}{\;错误\;}[]
    由于高频段的波长小，高频的衍射和绕射能力都强于低频。

\end{choice}


\begin{choice}{\;正确\;}[]
    通信系统的作用就是将信息从信源发送到一个或多个目的地。

\end{choice}


\begin{choice}{\;错误\;}[]
    5G 要求达到的峰值速率是 10Gbps。

\end{choice}


\begin{choice}{\;错误\;}[]
    C-V2X 只能采用 NR 实现，不能采用 LTE 技术实现。

\end{choice}

\begin{choice}{\;正确\;}[]
    5G 系统中 SRB3 是当 UE 处于 EN-DC 时使用 DCCH 逻辑信道的特殊的 RRC 消息。

\end{choice}


\begin{choice}{\;正确\;}[]
    智能家居属于 mMTC 典型应用。

\end{choice}


\begin{choice}{\;错误\;}[]
    5G 的基站功能重构为 CU 和 DU 两个功能实体，CU 主要处理物理层功能和实时性需求的层 2 功能。

\end{choice}


\begin{choice}{\;错误\;}[]
    大唐 5G 基站设备安装时，GPS/北斗馈线的长度大于 60 米，需要 1 处接地。

\end{choice}

\begin{choice}{\;正确\;}[]
    大唐5G设备安装时，电源线、地线与各种信号线缆水平间距大于150mm，不同设备的电源线禁止捆扎在一起：平行走线时，间距推荐大于100mm。

\end{choice}

\begin{choice}{\;错误\;}[]
    5G NR的PSS使用ZC序列，SSS使用m序列。

\end{choice}

\begin{choice}{\;错误\;}[]
    基站实现多流的空分复用（SDMA），相关性越强空分性能越好，可以提升单个用户的峰值速率。
\end{choice}

\begin{choice}{\;正确\;}[]
    5G控制面锚点不同区分为两大类：独立部署和非独立部署。

\end{choice}




